\documentclass{situationist-preprint}

\title{The Instability of Anarchy}
\msid{POL-07-a}
\msfamily{POL-07}
\msdiscipline{Political Philosophy}
\mstype{Principal manuscript}
\msstatus{Revised 7 October 2026}
\mskeywords{anarchism; state of nature; protective associations; institutions; persistence; Nozick}

\begin{document}
\maketitle

\begin{abstract}
Philosophical anarchism holds that no state is legitimate; political anarchism holds, further, that a society without a state is possible and desirable. We argue that a stateless society is coherent but unstable: it tends to develop institutions that perform the functions of a state. The argument has two parts. The first, adapted from Nozick, shows that protective associations in a stateless society tend to consolidate. The second, which is new, shows that institutions, once created, persist through periods of disuse, so that the consolidation is a ratchet: state-like institutions accumulate and are rarely dismantled.
\end{abstract}

\paragraph{Note on the revision.} Section~4 of the version of 6 October 2026 derived the ratchet effect from the companion manuscript POL-07-b, which has been withdrawn because it relied on an incorrect persistence condition for institutions. Section~4 now derives the ratchet effect directly from the corrected persistence condition of MET-20-a, Lemma~3.2. The main conclusion is unchanged.

\section{Anarchism}

Anarchism comes in two forms. \emph{Philosophical anarchism} holds that no existing state, and perhaps no possible state, has legitimate authority, so that there is no general duty to obey the law \citep{Wolff1970}. \emph{Political anarchism} holds in addition that a society without a state is possible and preferable, and that order can be maintained by voluntary cooperation \citep{Taylor1982}. Historical and anthropological work has shown that stateless societies exist and have existed for long periods \citep{Scott2009}, and that communities can manage common resources without central authority \citep{Ostrom1990}. The question is therefore not whether anarchy is possible but whether it is stable.

We take no position here on philosophical anarchism, which concerns legitimacy and is addressed in POL-02 and POL-03. Our claim concerns political anarchism: a stateless society is coherent, but there are forces that tend to turn it into a society with a state.

\section{Consolidation}

\citet{Nozick1974} argued that a minimal state could arise from a state of nature by an invisible-hand process that violates no one's rights. Individuals in a state of nature would form protective associations to enforce their rights. Since protection is a good whose value depends on the strength of the protector, and since conflicts between associations are costly, associations would tend to merge or to be absorbed by the strongest, until a dominant protective association emerged in each territory. This association, Nozick argued, would be a state in all but name.

\begin{proposition}\label{p:consolidate}
In a stateless society in which individuals form protective associations, the number of independent associations in a territory tends to decrease.
\end{proposition}

\begin{proof}[Sketch]
Protection exhibits increasing returns to scale: a larger association can deter more aggressors at lower cost per member, and in conflicts between associations the larger tends to prevail. Members of smaller associations therefore have reason to join larger ones, and smaller associations reason to merge. The process stops only when one association dominates, or when the costs of conflict lead associations to adopt common procedures for resolving disputes between their members, which is itself a step toward a common authority.
\end{proof}

The argument has been disputed. \citet{Leeson2014} argues that self-enforcing institutions in stateless settings can remain decentralized. We accept that consolidation is a tendency, not a law, and that it can be resisted by deliberate effort. The question is what happens when it is not resisted.

\section{The persistence of institutions}

MET-20-a, Lemma~3.2, establishes that an institution that has come into existence through collective acceptance persists as long as it is not collectively rejected by the community whose practices contain its rules. Active acceptance is required for an institution to come into being, but not for it to continue.

\section{The ratchet}

\begin{lemma}\label{l:ratchet}
State-like institutions created during consolidation persist through periods in which they are not used, unless they are collectively rejected.
\end{lemma}

\begin{proof}
A state-like institution created during consolidation, such as a common procedure for resolving disputes between members of different associations, comes into existence through collective acceptance by the members of the associations involved. By MET-20-a, Lemma~3.2, it persists until it is collectively rejected. A period in which the procedure is not invoked, because no dispute arises, is a period of disuse, not of rejection.
\end{proof}

\begin{theorem}\label{t:main}
A stateless society in which consolidation is not deliberately resisted tends to accumulate state-like institutions.
\end{theorem}

\begin{proof}
By Proposition~\ref{p:consolidate}, consolidation tends to create state-like institutions. By Lemma~\ref{l:ratchet}, those institutions persist unless collectively rejected. Collective rejection requires coordinated action by the community, which is costly and which, in the absence of a specific grievance, people have little reason to undertake. Creation is therefore more frequent than rejection, and the institutions accumulate.
\end{proof}

Theorem~\ref{t:main} explains a pattern that political anarchists themselves have noted: that stateless arrangements require constant vigilance, and that the institutions of coordination they create tend to outlive the occasions for which they were created. The theorem does not say that anarchy is impossible, only that its maintenance requires periodic collective rejection of institutions that have outlived their purpose, which is a demanding form of collective action.

\section{Objections}

\begin{objection}
The argument assumes that people do not reject unwanted institutions. But in a society committed to anarchist values, rejection of state-like institutions would be routine.
\end{objection}

\begin{reply}
We agree that a society with strong and widely shared anarchist commitments can resist the ratchet. The theorem concerns societies in which consolidation is not deliberately resisted. Its point is that the default, absent such commitment, runs toward the state, and that the commitment must be sustained, not merely held.
\end{reply}

\begin{objection}
Many stateless societies have lasted for centuries without developing states \citep{Scott2009}.
\end{objection}

\begin{reply}
Many of the societies Scott describes were stateless by design, in the sense that they actively resisted incorporation into neighbouring states and maintained practices, such as mobility and dispersed settlement, that made consolidation costly. They are cases of deliberate resistance, which the theorem allows. They do not show that consolidation does not occur where resistance is absent.
\end{reply}

\section{Conclusion}

Anarchy is coherent but unstable. Protective associations tend to consolidate, and the institutions consolidation creates persist through disuse unless they are collectively rejected. A stateless society that does not periodically reject such institutions will tend to acquire a state.

\begin{thebibliography}{99}\small
\bibitem[Leeson(2014)]{Leeson2014} Leeson, P.~T. (2014). \emph{Anarchy Unbound: Why Self-Governance Works Better Than You Think}. Cambridge University Press.
\bibitem[Nozick(1974)]{Nozick1974} Nozick, R. (1974). \emph{Anarchy, State, and Utopia}. Basic Books.
\bibitem[Ostrom(1990)]{Ostrom1990} Ostrom, E. (1990). \emph{Governing the Commons: The Evolution of Institutions for Collective Action}. Cambridge University Press.
\bibitem[Scott(2009)]{Scott2009} Scott, J.~C. (2009). \emph{The Art of Not Being Governed: An Anarchist History of Upland Southeast Asia}. Yale University Press.
\bibitem[Taylor(1982)]{Taylor1982} Taylor, M. (1982). \emph{Community, Anarchy and Liberty}. Cambridge University Press.
\bibitem[Wolff(1970)]{Wolff1970} Wolff, R.~P. (1970). \emph{In Defense of Anarchism}. Harper \& Row.
\end{thebibliography}

\end{document}
